% Slide 10 -- Genetic and multi-objective optimisation
\begin{frame}[plain]
\BlurBg
\SlideHeader{foundations}{Genetic and multi-objective optimisation}
\begin{tikzpicture}[remember picture,overlay]
  \ContentPanel

  % left: population -> selection -> offspring, with a small Pareto front
  \begin{scope}[shift={($(current page.south west)+(1.55cm,1.55cm)$)}]
    % population cloud
    \foreach \x/\y in {0.1/2.5,0.5/1.8,0.2/1.1,0.7/0.5,0.0/0.3,0.6/2.2,0.3/0.8}{%
      \fill[ink!28] (\x,\y) circle (1.6pt);}
    \node[text=mutedink,font=\fontsize{7}{8.5}\selectfont] at (0.35,-0.30) {population};
    \draw[-{Latex[length=1.8mm]},ink!45,line width=.5pt] (1.0,1.4) -- (2.0,1.4);
    % offspring (better)
    \foreach \x/\y in {2.5/2.0,2.8/1.3,2.4/0.9,2.9/0.5}{%
      \fill[terracotta!65] (\x,\y) circle (1.8pt);}
    \node[text=mutedink,font=\fontsize{7}{8.5}\selectfont] at (2.65,-0.30) {offspring};
    % pareto front on the right
    \draw[ink!45,line width=.4pt] (4.0,0) -- (5.9,0);
    \draw[ink!45,line width=.4pt] (4.0,0) -- (4.0,2.7);
    \draw[mintDeep!80!black,line width=.9pt]
      (4.15,2.5) .. controls (4.7,1.2) and (5.2,0.7) .. (5.85,0.45);
    \foreach \x/\y in {4.25/2.30,4.6/1.35,5.05/0.80,5.6/0.52}{%
      \fill[mintDeep!80!black] (\x,\y) circle (1.7pt);}
    \node[text=mutedink,font=\fontsize{7}{8.5}\selectfont] at (4.95,-0.30) {Pareto front};
  \end{scope}

  \node[anchor=north west,text=ink,font=\fontsize{9.4}{12.4}\selectfont,text width=6.30cm,align=left]
    at ($(current page.north west)+(8.10cm,-2.00cm)$) {%
    A genetic algorithm maintains a population of candidate solutions and
    improves it through selection, crossover, and mutation. Multi-objective
    variants balance competing goals along a Pareto front. The approach suits
    discrete design variables for which no gradient is available.};

  \PanelNote{Deciding which units to keep, under a fixed budget, is exactly this kind of discrete search.}
\end{tikzpicture}
\end{frame}
